\documentclass[11pt]{article}
\usepackage[margin=1in]{geometry}
\usepackage{amsmath,amssymb,amsthm,booktabs,hyperref}
\newtheorem{theorem}{Theorem}
\newtheorem{lemma}[theorem]{Lemma}
\newtheorem{proposition}[theorem]{Proposition}
\newtheorem{corollary}[theorem]{Corollary}
\newtheorem{conjecture}[theorem]{Conjecture}
\theoremstyle{definition}
\newtheorem{definition}[theorem]{Definition}
\newtheorem{remark}[theorem]{Remark}
\newtheorem{question}[theorem]{Question}
\newcommand{\E}{\mathbb{E}}
\renewcommand{\Re}{\operatorname{Re}}
\renewcommand{\Im}{\operatorname{Im}}
\title{Interference and the price of a feature:\\ what replaces Conjecture 20}
\author{}
\date{16 October 2026}
\begin{document}
\maketitle

\begin{abstract}
Conjecture 20 says that the unrestricted maximum $L_{2,b}$-margin of product networks equals the
interference-free (IF) value $\nu^{1-s}B_p\Gamma_s(p)$.  It is false: at $p=3$, $b=2$ certified networks beat it for
every $n=11,\dots,18$ (by a factor $2.21$ at $n=18$; see \texttt{mechanism\_note.pdf}), and for $n=12$ on
$b\in[2,2.13]$.  This note places the counterexamples inside the paper's framework.
(i) We write every margin condition in terms of one complex \emph{feature field} $\Phi_l:\mathbb Z_p^n\to\mathbb C$ per
frequency, and show that a single feature has margin $\ge c$ exactly when its field stays in the cone
$c+\{|\arg z|\le\pi/p\}$.
(ii) We define the \emph{price of a feature} $\pi_{n,b}(p)\in[\nu^{1-s},1]$, the best normalised margin of a single
feature with interference allowed, and prove the \emph{price--programme theorem}
$\gamma^*_b\ge\pi_{n,b}(p)B_p\Gamma_s(p)$, realised by networks whose amplitudes follow a maximiser $y^\star$
of the programme (so they use exactly $K^*(s,p)$ features).  Theorem~8 of the paper is the case
$\pi=\nu^{1-s}$; at $p=3$ the inequality is an equality.  Hence the counterexamples change the \emph{price of a feature},
not the spectral allocation.
(iii) We identify what the counterexamples exploit and prove it is unavailable to single-frequency networks at
$p\ge5$: there the feature field is \emph{circular} ($\E(\Phi-\E\Phi)^2=0$) and the cone has half-angle $\pi/p<\pi/4$,
which forces a margin loss $\ge k_p\sigma$ proportional to the interference $\sigma$.  At $p=3$ the field need not be
circular because $2\equiv-1\pmod 3$, and the certified $n=11$ network uses exactly this freedom: interference
$258$ times the feature's energy, with imaginary fluctuations carrying $1.78$ times the energy of the real ones, parked
on the $60^\circ$ cone edges.  Whether the loss always outweighs the saving at $p\ge5$ reduces to one constant
$k^*(n,s)$ (estimated numerically).
(iv) We list what in the paper is untouched, and give the corrected statement of Appendix~D.
\end{abstract}

\section{Setting}
$p$ is an odd prime, $\rho=e^{2\pi i/p}$, $F=\{1,\dots,(p-1)/2\}$, $\nu=n+1$, $2\le b\le\nu$, $s=\nu/b$, $r=1/s$,
$\kappa_l(\delta)=1-\cos(2\pi l\delta/p)$ and $\kappa_{\min}(p)=1-\cos(2\pi/p)=\min_{\delta\neq0}\kappa_l(\delta)$.
A product network is $f(x)[c]=\sum_i w_i(c)\prod_{t=1}^n u_{it}(x_t)$, its margin is
$h=\min_x\min_{c\ne y(x)}(f(x)[y(x)]-f(x)[c])$, $y(x)=\sum_t x_t$, and $\gamma=h/\|\theta\|^\nu_{2,b}$; $\gamma^*_b$ is the
supremum over all widths.  $B_p=2(2p\nu)^{-\nu/2}$ and $C_b^{\nu/b}=B_p$ is the single-neuron budget
(paper, Appendix~C): $\sum_{j}|\beta_{ij}|^{r}\le C_b\|\omega_i\|^b$.  $\beta_l$, $\beta_{l,T}$ are the paper's clean and
sign-flipped coefficients (Section~4), and $\beta_{l,T}=\sum_i\beta_{il}\xi_{il}^{T}$ with $|\xi_{ilt}|=1$
(paper, Lemma~6).  The interference-free (IF) value is $\gamma^*_{b,\mathrm{IF}}=\nu^{1-s}B_p\Gamma_s(p)$ (Theorem~8).

\section{Feature fields and the cone}
\begin{definition}
For a network $\theta$ and $l\in F$ let
$\Phi_l(x):=\frac2p\,\rho^{l\,y(x)}\sum_{c}f(x)[c]\rho^{-lc}$.
\end{definition}

\begin{lemma}[decomposition]\label{lem:dec}
$f(x)[c]=\bar f(x)+\sum_{l\in F}\Re\big(\Phi_l(x)\rho^{l(c-y(x))}\big)$, hence for every $x$ and $\delta\ne0$
\[
g(x,\delta):=f(x)[y(x)]-f(x)[y(x)-\delta]=\sum_{l\in F}\Re\big(\Phi_l(x)(1-\rho^{-l\delta})\big).
\]
Moreover $\E_x\Phi_l=\overline{\beta_l}$, and for a product network
$\Phi_l(x)=\frac2p\sum_i\hat w_i(l)\prod_t u_{it}(x_t)\rho^{lx_t}$.
\end{lemma}
\begin{proof}
Fourier inversion of the real function $c\mapsto f(x)[c]$ gives
$f(x)[c]=\bar f(x)+\sum_{l\in F}\frac2p\Re(\hat f_x(l)\rho^{lc})$ with $\hat f_x(l)=\sum_cf(x)[c]\rho^{-lc}$, and
$\frac2p\hat f_x(l)\rho^{lc}=\Phi_l(x)\rho^{l(c-y)}$.  For the mean,
\[\beta_l=2p^{-\nu}\sum_{x,c}f(x)[c]\rho^{-l\sum_tx_t+lc}=2p^{-\nu}\sum_x\overline{\rho^{ly(x)}\hat f_x(l)}=\overline{\E_x\Phi_l},\]
using $\rho^{l\sum_t x_t}=\rho^{ly(x)}$.  The last formula is $\rho^{ly(x)}=\prod_t\rho^{lx_t}$.
\end{proof}

\begin{lemma}[cone]\label{lem:cone}
Let $z\in\mathbb C$, $c\in\mathbb R$, $l\in F$, and $K_p=\{w:|\Im w|\le\tan(\pi/p)\Re w\}$.  Then
$\Re(z(1-\rho^{-l\delta}))\ge c\,\kappa_l(\delta)$ for all $\delta\ne0$ if and only if $z\in c+K_p$, i.e.\
$\mathrm{cm}(z):=\Re z-\cot(\pi/p)|\Im z|\ \ge c$.  If $\mathrm{cm}(z)\ge0$ then
$\min_{\delta\ne0}\Re(z(1-\rho^{-l\delta}))=\kappa_{\min}(p)\,\mathrm{cm}(z)$.
\end{lemma}
\begin{proof}
Write $z=a+ib$ and $\vartheta=2\pi l\delta/p$; as $\delta$ runs over $\mathbb Z_p\setminus\{0\}$, $\vartheta/2$ runs over
$\{\pi k/p:1\le k\le p-1\}$.  Since $1-e^{-i\vartheta}=2\sin(\vartheta/2)\,(\sin(\vartheta/2)+i\cos(\vartheta/2))$,
\[\Re(z(1-e^{-i\vartheta}))=2\sin(\vartheta/2)\big(a\sin(\vartheta/2)-b\cos(\vartheta/2)\big),\qquad
\kappa=2\sin^2(\vartheta/2),\]
and $\sin(\vartheta/2)>0$, so the inequality reads $a-c\ge b\cot(\vartheta/2)$.  This is affine in
$\cot(\vartheta/2)\in[-\cot(\pi/p),\cot(\pi/p)]$, whose extreme values occur at $k=1,p-1$; so it holds for all $k$ iff
$a-c\ge|b|\cot(\pi/p)$.  For the last claim take $c=\mathrm{cm}(z)$: every gap is $\ge c\kappa_l(\delta)\ge c\kappa_{\min}$,
and the extreme $k\in\{1,p-1\}$ with $b\cot(\vartheta/2)=|b|\cot(\pi/p)$ has $\kappa=\kappa_{\min}$ and equality.
\end{proof}

For a single feature the margin is therefore governed by the cone $c+K_p$ of half-angle $\pi/p$: $60^\circ$ at $p=3$,
$36^\circ$ at $p=5$, $25.7^\circ$ at $p=7$.

\section{The price of a feature and the price--programme theorem}
\begin{definition}
A network is \emph{single-feature} if every output vector $w_i$ lies in
$\mathrm{span}\{\mathbf 1,\cos(2\pi\cdot/p),\sin(2\pi\cdot/p)\}$; then only $\Phi_1$ is non-zero.  The \emph{price of a
feature} is
\[
\pi_{n,b}(p):=\sup\Big\{\frac{\min_x\mathrm{cm}(\Phi_1(x))}{B_p\,\|\theta\|_{2,b}^{\nu}}:\ \theta\text{ single-feature,
any width}\Big\}.
\]
\end{definition}
By Lemma~\ref{lem:cone}, the largest normalised margin of a single-feature network is
$\kappa_{\min}(p)\,B_p\,\pi_{n,b}(p)$.

\begin{proposition}\label{prop:range}
$\nu^{1-s}\le\pi_{n,b}(p)\le1$.
\end{proposition}
\begin{proof}
Lower bound: the paper's block (Proposition~9) at $j=1$ has $\Phi_1\equiv A>0$ and attains Theorem~8 with $y=e_1$, i.e.\
normalised margin $\nu^{1-s}B_p\kappa_{\min}$.  Upper bound: $\mathrm{cm}(\Phi_1(x))\le\Re\Phi_1(x)$, so
$\min_x\mathrm{cm}\le\Re\E\Phi_1=\Re\overline{\beta_1}\le\sum_i|\beta_{i1}|\le(\sum_i|\beta_{i1}|^r)^{1/r}
\le(C_b\|\theta\|^b_{2,b})^{s}=B_p\|\theta\|_{2,b}^\nu$, using $r\le1$.
\end{proof}

\begin{theorem}[price--programme]\label{thm:pp}
For every odd prime $p$, $n\ge1$, $2\le b\le\nu$ and every $y\in\mathbb R^F_{\ge0}$ with $\sum_ly_l^{1/s}\le1$,
\[
\gamma^*_b\ \ge\ \pi_{n,b}(p)\,B_p\,\min_{\delta\ne0}\sum_{l}y_l\kappa_l(\delta),\qquad\text{in particular}\qquad
\gamma^*_b\ \ge\ \pi_{n,b}(p)\,B_p\,\Gamma_s(p).
\]
The networks realising the bound are frequency-relabelled copies of one single-feature network, one copy per
$l\in\operatorname{spt}(y)$, scaled by $y_l$; at $y=y^\star$ they use exactly $K^*(s,p)$ features with amplitudes
proportional to $y^\star$.
\end{theorem}
\begin{proof}
Let $\theta$ be single-feature with $c:=\min_x\mathrm{cm}(\Phi_1(x))>0$ and $P:=\|\theta\|^b_{2,b}$.  For $l\in F$ define
$\theta^{(l)}$ by $u^{(l)}_{it}(a)=u_{it}(la)$, $w^{(l)}_i(c')=w_i(lc')$.  Since $a\mapsto la$ permutes $\mathbb Z_p$,
$\|\theta^{(l)}\|_{2,b}=\|\theta\|_{2,b}$, and $f^{(l)}(x)[c']=f(lx)[lc']$ with $y(lx)=l\,y(x)$, so
$g^{(l)}(x,\delta)=g(lx,l\delta)\ge c\,\kappa_1(l\delta)=c\,\kappa_l(\delta)$ by Lemma~\ref{lem:cone}.
Scaling every weight of $\theta^{(l)}$ by $y_l^{1/\nu}$ multiplies its table by $y_l$ and its $\|\cdot\|^b_{2,b}$ by
$y_l^{r}$.  The union $\Theta$ of the scaled copies has $g_\Theta(x,\delta)=\sum_ly_lg^{(l)}(x,\delta)\ge
c\sum_ly_l\kappa_l(\delta)$ and $\|\Theta\|^b_{2,b}=P\sum_ly_l^r\le P$, so
$\gamma(\Theta)\ge c\min_\delta\sum_ly_l\kappa_l(\delta)/P^{s}$.  Taking the supremum over $\theta$ gives the claim.
\end{proof}

\begin{corollary}[$p=3$]\label{cor:p3}
$F=\{1\}$, so every network is single-feature and $\gamma^*_b(3)=\tfrac32\,B_3\,\pi_{n,b}(3)$.  The certified
networks give $\pi_{n,2}(3)\,\nu^{s-1}\ge1.0630$, $1.1933$, $1.1836$, $1.4568$, $1.6170$, $1.5954$, $2.0110$, $2.2119$ for
$n=11,\dots,18$ (one-dimensional networks of \texttt{mechanism\_note.pdf}; the first notes had $1.0543$, $1.1091$,
$1.4293$ for $n=11,13,14$), and $\pi_{12,b}(3)>13^{1-13/b}$ for $b\in[2,2.13]$,
$\pi_{11,b}(3)>12^{1-12/b}$ for $b\in[2,2.02]$.
\end{corollary}

\begin{corollary}[failure is a statement about the price]\label{cor:transfer}
Conjecture~20 holds at $(p,n,b)$ only if $\pi_{n,b}(p)=\nu^{1-s}$.  Whenever a single feature can be bought below the
interference-free price ($\pi_{n,b}(p)>\nu^{1-s}$), the unrestricted optimum exceeds $\gamma^*_{b,\mathrm{IF}}$ by at least the
factor $\pi_{n,b}(p)/\nu^{1-s}$ \emph{with the same spectral allocation $y^\star$ and the same number $K^*(s,p)$ of
features}.
\end{corollary}

\begin{remark}[what the counterexamples do not touch]
At $p=3$ there is a single feature, so the counterexamples change only the constant in front of $B_3\Gamma_s(3)$.
Theorem~4 (universal bound, i.e.\ $\pi\le1$), Theorem~8 (exact IF value and IF structure), Theorem~10 (order
$p^{-\nu/2}(\ln p)^{1-s}$ of the unrestricted optimum), Theorem~11 (programme maximisers use $\Theta(\log p)$ features) and
Corollary~12 (unrestricted near-maximisers need $\Omega(\log p)$ features) are statements that the counterexamples do not
contradict.  The smallest known counterexample has $\nu=12$; the paper's experiments use $\nu\le5$ and $p\ge11$.
\end{remark}

\section{Why $p=3$: circular feature fields}
\begin{definition}
A network is \emph{single-frequency} (at frequency $1$) if every $u_{it}$ and every $w_i$ lies in
$\mathrm{span}\{\mathbf 1,\cos(2\pi\cdot/p),\sin(2\pi\cdot/p)\}$: constant components are allowed.  At $p=3$ every
network is single-frequency.  Relabelling $a\mapsto la$ (as in the proof of Theorem~\ref{thm:pp}) transfers everything
below to a network using any single frequency $l$.
\end{definition}

\begin{lemma}[support]\label{lem:supp}
For a single-frequency network, $\Phi_1(x)=\sum_{e\in\{0,1,2\}^n}a_e\rho^{e\cdot x}$ with $a_0=\overline{\beta_1}$ and
$a_{2\cdot\mathbf 1_T}=\overline{\beta_{1,T}}$ for every $T\subseteq[n]$; the coefficients with some $e_t=1$ come from
constant input components.
\end{lemma}
\begin{proof}
$u(a)\rho^{a}=p^{-1}\big(\hat u(-1)+\hat u(0)\rho^{a}+\hat u(1)\rho^{2a}\big)$; expand the product in
Lemma~\ref{lem:dec}.  The term taking $\hat u_{it}(1)$ for $t\in T$ and $\hat u_{it}(-1)$ otherwise equals
$\frac2pp^{-n}\sum_i\hat w_i(1)\prod_{t\in T}\hat u_{it}(1)\prod_{t\notin T}\hat u_{it}(-1)
=\overline{2p^{-\nu}\hat f_\theta(k;-1)}=\overline{\beta_{1,T}}$ with $k_t=-1$ on $T$ and $+1$ off $T$.
\end{proof}

\begin{lemma}[circularity]\label{lem:circ}
If $p\ge5$, every single-frequency network satisfies $\E_x(\Phi_1-\E\Phi_1)^2=0$.  At $p=3$ this fails in general.
\end{lemma}
\begin{proof}
$\E_x\rho^{(e+e')\cdot x}=1$ iff $e_t+e'_t\equiv0\pmod p$ for all $t$, and $0$ otherwise.  For
$e,e'\in\{0,1,2\}^n$ we have $e_t+e'_t\in\{0,\dots,4\}$, which for $p\ge5$ is $\equiv0$ only if $e_t=e'_t=0$.  So
$\E(\Phi_1-a_0)^2=\sum_{e,e'\ne0}a_ea_{e'}\E\rho^{(e+e')\cdot x}=0$.  At $p=3$, $1+2\equiv0$: the constant-input
characters and the sign-flipped characters are complex conjugates ($\omega^{2x}=\overline{\omega^{x}}$), and the
pseudo-variance is unconstrained.
\end{proof}

\begin{lemma}[circular cone bound]\label{lem:cc}
Let $Z$ be a complex random variable on a finite probability space with $\E Z=A$, $\E(Z-A)^2=0$ and
$Z\in c+K_\alpha$ almost surely, where $K_\alpha=\{|\Im w|\le\tan\alpha\,\Re w\}$ and $0<\alpha<\pi/4$.  Then
\[\Re A-c\ \ge\ \sqrt{\cot^2\alpha-1}\;\sigma,\qquad \sigma^2:=\tfrac12\E|Z-A|^2 .\]
Without circularity, for any $Z$ with $Z\in c+K_\alpha$,
$(1-t^2)\E|Z-A|^2\le 2t^2(\Re A-c)^2+(1+t^2)\Re\E(Z-A)^2$, $t=\tan\alpha$.
\end{lemma}
\begin{proof}
Put $X=\Re Z-c\ge0$, $Y=\Im Z$, so $|Y|\le tX$.  Let $E=\E|Z-A|^2$ and $P=\Re\E(Z-A)^2$; then
$\mathrm{Var}X=(E+P)/2$ and $\mathrm{Var}Y=(E-P)/2$.  Hence
$(E-P)/2\le\E Y^2\le t^2\E X^2=t^2\big((\Re A-c)^2+(E+P)/2\big)$, which is the second claim; with $P=0$ and
$E=2\sigma^2$ it gives $\sigma^2(1-t^2)\le t^2(\Re A-c)^2$, the first.
\end{proof}

\begin{theorem}[no $p=3$ mechanism for $p\ge5$]\label{thm:circ}
Let $p\ge5$ and $\theta$ be single-frequency.  With $k_p:=\sqrt{\cot^2(\pi/p)-1}$ and
$\sigma^2:=\frac12\E|\Phi_1-\E\Phi_1|^2\ \ge\ \frac12\sum_{\emptyset\ne T\subseteq[n]}|\beta_{1,T}|^2$,
\[
h(\theta)\ \le\ \kappa_{\min}(p)\,\big(\Re\overline{\beta_1}-k_p\,\sigma\big),
\qquad
\gamma(\theta)\ \le\ \kappa_{\min}(p)\,B_p\,\Lambda_{n,s}(k_p),
\]
where $\Lambda_{n,s}(k):=\sup\big\{(\Re A-k\sqrt{E/2})/(\sum_i|\beta_i|^{r})^{s}\big\}$ over finitely many atoms
$\beta_i\in\mathbb C$, $\xi_i\in\mathbb T^n$, with $A=\sum_i\beta_i$, $E=\sum_{\emptyset\ne T}|\sum_i\beta_i\xi_i^T|^2$.
$\Lambda_{n,s}$ is non-increasing, $\Lambda_{n,s}(k)\ge\nu^{1-s}$, and $\Lambda_{n,s}(k)=\nu^{1-s}$ for all
$k\ge k^*(n,s):=\sup(\Re A-\nu^{1-s}(\sum|\beta_i|^r)^s)/\sqrt{E/2}$.  Consequently, if $k_p\ge k^*(n,s)$, no
single-frequency network at $(p,n,b)$ beats the IF value of a single feature.
\end{theorem}
\begin{proof}
By Lemma~\ref{lem:cone} (with $h>0$, else nothing to prove), $h=\kappa_{\min}\min_x\mathrm{cm}(\Phi_1(x))$, and
$\Phi_1\in c+K_{\pi/p}$ with $c=\min_x\mathrm{cm}$.  Lemma~\ref{lem:circ} and Lemma~\ref{lem:cc} with $\alpha=\pi/p<\pi/4$
give $c\le\Re\E\Phi_1-k_p\sigma$, and Parseval with Lemma~\ref{lem:supp} gives
$2\sigma^2=\sum_{e\ne0}|a_e|^2\ge\sum_{T\ne\emptyset}|\beta_{1,T}|^2$.  For the second bound, divide by
$\|\theta\|^\nu_{2,b}\ge(C_b^{-1}\sum_i|\beta_{i1}|^r)^s=B_p^{-1}(\sum_i|\beta_{i1}|^r)^s$ and use
$\beta_{1,T}=\sum_i\beta_{i1}\xi_{i}^{T}$ (paper, Lemma~6) and $\Re\overline{\beta_1}=\Re\sum_i\beta_{i1}$.
The IF block gives $\Lambda\ge\nu^{1-s}$.  Atoms with $E=0$ satisfy $\Re A\le\nu^{1-s}(\sum|\beta_i|^r)^s$ by the paper's
Lemma~7, and for atoms with $E>0$ the inequality $(\Re A-k\sqrt{E/2})\le\nu^{1-s}(\sum|\beta_i|^r)^s$ for all
$k\ge k^*$ is the definition of $k^*$.
\end{proof}

\begin{proposition}[matrix reduction]\label{prop:mat}
$k^*(n,s)\le\sqrt2\,\kappa_M(\nu,s)$, where
\[\kappa_M(\nu,s):=\sup\Big\{\frac{1-\nu^{-s}\|H\|_{S_{r}}}{\|\mathrm{strictlower}(H)\|_F}:\ H\in\mathbb C^{\nu\times\nu},\ H_{jj}=1\Big\}\]
and $\|H\|_{S_r}=(\sum_k\sigma_k(H)^r)^{1/r}$.
\end{proposition}
\begin{proof}
Rotate the atoms so that $A>0$ (this does not change $E$ or the cost, and $\Re A\le|A|$).  Let
$m_i=(1,\xi_{i1},\xi_{i1}\xi_{i2},\dots,\xi_{i1}\cdots\xi_{in})$ and $H=\sum_i\beta_im_im_i^*$.  Its diagonal is $A$,
and its strictly lower entries are the interval sums $\sum_i\beta_i\xi_i^{(k,j]}$, distinct non-empty patterns, so
$\|\mathrm{strictlower}H\|_F^2\le E$.  As $\|\beta_im_im_i^*\|_{S_r}=\nu|\beta_i|$ and
$\|X+Y\|_{S_r}^r\le\|X\|_{S_r}^r+\|Y\|_{S_r}^r$ for $0<r\le1$ (McCarthy, Rotfel'd),
$\sum_i|\beta_i|^r\ge\nu^{-r}\|H\|_{S_r}^r$, i.e.\ $\nu^{1-s}(\sum|\beta_i|^r)^s\ge\nu^{-s}\|H\|_{S_r}$.  Hence
$A-\nu^{1-s}(\sum|\beta_i|^r)^s\le A(1-\nu^{-s}\|H/A\|_{S_r})\le\kappa_M\|\mathrm{strictlower}H\|_F\le\sqrt2\kappa_M\sqrt{E/2}$.
\end{proof}

\section{Anatomy of the counterexample}
For the certified $p=3$, $n=11$ network (normalised so that $\mathrm{cm}$ is the ratio to the IF value), the field
$\Phi=\Phi_1$ satisfies:
\begin{center}
\begin{tabular}{ll}
\toprule
clean amplitude $A=\E\Phi$ & $2.3960+0.0138i$ \\
cone margin $c=\min\mathrm{cm}(\Phi)$ & $1.05440=0.440\,\Re A$ \\
interference energy $\E|\Phi-A|^2$ & $258.4\,|A|^2$ \\
pseudo-variance $\E(\Phi-A)^2$ & $(-72.2+10.0i)\,|A|^2$ \\
$\mathrm{Var}\,\Im\Phi/\mathrm{Var}\,\Re\Phi$ & $1.776$ \\
$\max_x|\Phi(x)|$ & $4491\,|A|$ \\
count classes with $|\arg(\Phi-c)|>59.9^\circ$ & $21$ of $78$ \\
\bottomrule
\end{tabular}
\end{center}
The network gives up $56\%$ of its clean amplitude and carries interference two orders of magnitude larger than the
feature, but the interference is aligned with the two edges of the $60^\circ$ cone: on those inputs the correct logit
ties (up to the margin) with one neighbouring class and the third class is pushed far down.  This is possible only
because the fluctuation is \emph{non-circular} ($\mathrm{Var}\,\Im>\mathrm{Var}\,\Re$, allowed since $\tan60^\circ>1$),
and that is possible only because $2\equiv-1\pmod3$ (Lemma~\ref{lem:circ}).  In exchange it uses $8$ neurons instead of
$\nu=12$, and $(12/8)^{s-1}$-type savings in the concave $\ell_r$ cost outweigh the lost amplitude.

\section{Numerical evidence beyond $p=3$}
All numbers in this section come from local optimisation (L-BFGS with an annealed soft-min of the margins) and are
\emph{evidence}, not certificates; scripts are in the supplementary folder.

\paragraph{1. The robust constant $k^*(n,s)$.}  Lower bounds on the supremum in Theorem~\ref{thm:circ}, at $b=2$
($s=\nu/2$), from optimising over atoms (general chains for $n\le7$, slot-symmetric chains above):
\begin{center}
\begin{tabular}{lccccccc}
\toprule
$n$ & 2 & 3 & 5 & 7 & 9 & 11 & 13\\
$k^*(n,\nu/2)\ \ge$ & 0.345 & 0.543 & 0.764 & 0.928 & 1.055 & 1.073 & 1.099\\
\bottomrule
\end{tabular}
\qquad
\begin{tabular}{lcccc}
\toprule
$p$ & 5 & 7 & 11 & 13\\
$k_p$ & 0.946 & 1.820 & 3.256 & 3.932\\
\bottomrule
\end{tabular}
\end{center}
The growth flattens.  If these lower bounds are close to the supremum, Theorem~\ref{thm:circ} excludes beating the IF value
with single-frequency networks for $p\ge7$ and all $n\le13$, and for $p=5$ and $n\le7$.  For the matrix constant of
Proposition~\ref{prop:mat} we found $\kappa_M(3,1.5)\approx0.289$, $\kappa_M(4,2)\approx0.476$, $\kappa_M(6,3)\approx0.694$.

\paragraph{2. The $L^2$ bound is far from tight pointwise.}  At $p=5$, $n=9$ the best atom configuration for
$(\Re A-k_5\sigma)/\nu^{1-s}C^s$ reaches $1.186$ ($m=7$ atoms, $E=0.79|A|^2$).  Realised as a single-frequency network,
its cone efficiency is $-31$, and local optimisation from it (constant input components allowed) never reaches a positive
margin.  The pointwise cone constraint is much stronger than its second moment.

\paragraph{3. IF behaves as a local maximum.}  Starting from the IF block plus noise, with the cone widened to any
half-angle between $36^\circ$ and $89^\circ$, local optimisation in the full single-feature class (arbitrary input
vectors) returns to the IF value ($0.9999$) at $p=3$ and $p=5$, $n=11$.  The $p=3$ counterexamples live in a different
basin, reached only through a structured family of $\approx0.7\nu$ neurons.

\paragraph{4. Narrowing the cone kills the counterexample.}  The certified $p=3$, $n=11$ network evaluated with a
half-angle of $59^\circ$ instead of $60^\circ$ has cone efficiency $-218$ (and $-2438$ at $50^\circ$), because its
interference sits on the cone edges.

\paragraph{5. Fewer than $\nu$ neurons.}  The counterexamples save cost by using fewer than $\nu$ neurons for one feature.  Best normalised margin (ratio
to the IF value) over 24 random restarts of slot-symmetric networks with exactly $m=\nu-1$ neurons:
\begin{center}
\begin{tabular}{lcccccccccc}
\toprule
$(p,n)$ & $(3,3)$ & $(3,4)$ & $(5,2)$ & $(5,3)$ & $(5,4)$ & $(5,5)$ & $(5,6)$ & $(5,7)$ & $(7,2)$ & $(7,3)$\\
ratio & 0.599 & 0.597 & $\le0$ & $\le0$ & 0.278 & 0.310 & $\le0$ & $\le0$ & $\le0$ & $\le0$\\
\bottomrule
\end{tabular}
\end{center}
At $p=3$, networks with at most $\nu-1$ neurons classify at every tested $n$ ($2\le n\le14$), keep $\approx60\%$ of the
IF value already at $n=3,4$, and form the branch that crosses $1$ at $n=11$.  At $p=5$ such networks were found only at
$n=4,5$, at $28$--$31\%$ of the IF value; at $p=7$ none were found for $n\le5$ (a continuation run from $n=5$ to $14$
also found none).  Random restarts are weak at larger $n$ (at $p=3$ they too miss the structured branch beyond $n\approx6$),
so this is evidence, not proof.


\paragraph{6. Cross-feature compensation.}  Conjecture~\ref{conj:pp} can only fail through features that use each
other's slack.  We started from the IF network with the programme allocation $y^\star$ (several features), added noise,
optionally deleted neurons, and optimised the true margin over all $p^n$ inputs (general, not slot-symmetric networks):
\begin{center}
\begin{tabular}{lcccc}
\toprule
$(p,n)$, features $K^*$, IF width & from IF $+$ noise & one neuron deleted & two deleted\\
\midrule
$(5,2)$, $K^*=2$, $6$ & $1.000000$ & $0.837$ & $0.782$\\
$(7,2)$, $K^*=3$, $9$ & $1.000000$ & $0.980$ & --\\
$(7,3)$, $K^*=2$, $8$ & $0.999995$ & $0.764$ & --\\
\bottomrule
\end{tabular}
\end{center}
(Ratios to the IF value; best of 3--4 seeds.)  In every case the optimiser returns to the IF network, and removing
neurons costs margin; $0.782=\kappa_{\min}(5)/\Gamma_{3/2}(5)$ is the single-feature IF network.  Random restarts at
$(5,2)$ with widths $4$--$8$ (30 runs) never exceeded $0.782$.  No cross-feature gain was found.



\section{Corrected statement of Appendix D}
\begin{theorem}[replaces Conjecture 20]
For all odd primes $p$, $n\ge1$ and $2\le b\le\nu$,
\[
\nu^{1-s}B_p\Gamma_s(p)\ \le\ \pi_{n,b}(p)\,B_p\,\Gamma_s(p)\ \le\ \gamma^*_b\ \le\ B_p\Gamma_s(p),
\]
with equality in the middle at $p=3$.  The first inequality is strict at $p=3$, $b=2$, $n\in\{11,13,14\}$ and on the
$b$-intervals of Corollary~\ref{cor:p3}; it is an equality for $b_0\le b\le\nu$ (previous notes, Part~II) and at
$p=3$, $n=2$ for all $b$ (Part~I).
\end{theorem}

\begin{conjecture}[price--programme equality]\label{conj:pp}
$\gamma^*_b=\pi_{n,b}(p)\,B_p\,\Gamma_s(p)$: interference changes the price of a feature but not the programme.
\end{conjecture}

\begin{question}[is $p=3$ the only exception?]\label{q:p5}
Is $\pi_{n,b}(p)=\nu^{1-s}$ for every $p\ge5$?  Theorem~\ref{thm:circ} gives a positive answer for single-frequency
networks whenever $k_p\ge k^*(n,s)$.  Open are (a) single-frequency networks with $k_p<k^*(n,s)$ (numerically
$p=5$, $n\ge9$), and (b) networks whose neurons carry the partner frequencies $2l$, $3l$ of a feature $l$: these are the
only way to make a feature field non-circular at $p\ge5$, and by the second part of Lemma~\ref{lem:cc} a large
interference needs a large positive pseudo-variance.
\end{question}

\begin{remark}[update]
The companion \texttt{mechanism\_note.pdf} identifies the $p=3$ mechanism exactly (shared sign pattern, field
$e^{i\pi\tau/p}\sum_iV_ie^{\lambda_i\tau}$), proves that such networks need at least $\nu-\lfloor\nu/p\rfloor$ neurons (so
none are saved when $p>\nu$), finds the mechanism losing to IF at $p=5,7$ by a factor growing in $n$, and exhibits genuine
cross-feature compensation at $p=5$ (an obstacle to any feature-by-feature proof of Conjecture~\ref{conj:pp}).
\end{remark}

\paragraph{Suggested wording for the paper.}  Section~4 (``How many features''): ``Whether a network carrying
interference can beat the class value is an open conjecture'' $\to$ ``Interference can beat the class value: at $p=3$
and $\nu\ge12$ a single feature can be bought below the interference-free price (Appendix~D).  This changes the price of
a feature, not the allocation: the best networks built from cheaper features still follow the programme
(Theorem~D.2).''  Limitations: ``The exact optimum is proved for interference-free networks; for all networks we prove
$\pi_{n,b}(p)B_p\Gamma_s(p)\le\gamma^*_b\le B_p\Gamma_s(p)$, where $\pi_{n,b}(p)$, the price of a feature, equals the
interference-free price $\nu^{1-s}$ near $b=\nu$ and exceeds it at $p=3$, $\nu\ge12$, $b=2$.''

\section{Open problems}
\begin{enumerate}
\item Conjecture~\ref{conj:pp}.  The obstacle is cross-feature compensation: a feature field may leave its own cone
where another feature has slack, and at $b=2$ one neuron can feed several features at no budget cost.
\item Question~\ref{q:p5}, in particular single-frequency networks at $p=5$ with $n\ge9$, where the numerical value of
$k^*(n,s)$ exceeds $k_5=0.946$, and partner-frequency networks at every $p\ge5$.
\item The growth of $\pi_{n,2}(3)\nu^{s-1}$ in $n$: certified $1.06$ at $n=11$ up to $2.21$ at $n=18$, about $10\%$ per
unit of $n$ (\texttt{mechanism\_note.pdf}, which also proves that the shared-sign networks behind these values need at
least $\nu-\lfloor\nu/p\rfloor$ neurons).
\end{enumerate}
\end{document}
